% ============================================================
%  SCHOLA DOMESTICA — Magister Mathematica
%  Level 1 | Week 16 | Day 2
%  Topic: Roman Numerals I–X Review (Enrichment) + Oral Drill #4
% ============================================================
\documentclass[12pt, letterpaper]{article}
\usepackage{sd_style}

\renewcommand{\lessonweek}{16}
\renewcommand{\lessonnum}{2}
\renewcommand{\lessontitle}{Roman Numerals I–X \& Oral Drill \#4}

% IMAGE PROMPT (Beatrix Potter watercolour style):
% A wise old tortoise examines a stone tablet carved with Roman numerals I through X
% in a sunny garden. Climbing roses on an old wall behind. White background, no text.

\begin{document}

\lessonheader{Week 16}{2}{Roman Numerals I--X \& Oral Drill \#4}

{\small Name:\ansblanklg \hfill Date:\ansblank}
\goldrule

\begin{mdframed}[backgroundcolor=sd-gold!15,linecolor=sd-gold,linewidth=1.5pt,roundcorner=4pt,
  innerleftmargin=10pt,innerrightmargin=10pt,innertopmargin=8pt,innerbottommargin=8pt]
\textbf{\color{sd-gold}Oral Speed Drill \#4 — Final Unit B Drill}\\[4pt]
\textit{Complete mixed facts. Goal: 22 of 25.}\\[6pt]
\begin{multicols}{5}
\begin{enumerate}[leftmargin=*,label=\small\arabic*.,itemsep=1pt]
  \item $6+4$ \textit{[10]}  \item $10-4$ \textit{[6]}   \item $7+2$ \textit{[9]}
  \item $9-2$ \textit{[7]}   \item $5+5$ \textit{[10]}   \item $10-5$ \textit{[5]}
  \item $8+1$ \textit{[9]}   \item $9-1$ \textit{[8]}    \item $4+3$ \textit{[7]}
  \item $7-3$ \textit{[4]}   \item $5+3$ \textit{[8]}    \item $8-3$ \textit{[5]}
  \item $6+3$ \textit{[9]}   \item $9-3$ \textit{[6]}    \item $4+6$ \textit{[10]}
  \item $10-6$ \textit{[4]}  \item $3+7$ \textit{[10]}   \item $10-7$ \textit{[3]}
  \item $2+8$ \textit{[10]}  \item $10-8$ \textit{[2]}   \item $9+1$ \textit{[10]}
  \item $10-9$ \textit{[1]}  \item $4+5$ \textit{[9]}    \item $9-5$ \textit{[4]}
  \item $6+2$ \textit{[8]}
\end{enumerate}
\end{multicols}
\noindent Score:\underline{\hspace{1.5cm}}/25
\end{mdframed}

\bigskip
\noindent{\color{sd-blue}\large\textbf{Enrichment — Roman Numerals I–X}} \hfill \textit{(Mastery Review)}

\bigskip
\begin{instructbox}
\begin{center}
\begin{tabular}{*{10}{c}}
\textbf{I} & \textbf{II} & \textbf{III} & \textbf{IV} & \textbf{V} &
\textbf{VI} & \textbf{VII} & \textbf{VIII} & \textbf{IX} & \textbf{X} \\[4pt]
1 & 2 & 3 & 4 & 5 & 6 & 7 & 8 & 9 & 10
\end{tabular}
\end{center}
\medskip
\textbf{Rules:} V = 5 \quad I before V means subtract: IV = 4 \quad I after V means add: VI = 6
\end{instructbox}

\bigskip
\noindent{\color{sd-blue}\textbf{Part I — Write the Arabic numeral.}} \hfill \textit{(Problems 1–8)}

\bigskip
\noindent\begin{tabular}{@{}p{3cm}p{3cm}p{3cm}p{3cm}@{}}
\prob III $=\ansblank$ & \prob VII $=\ansblank$ & \prob IV $=\ansblank$ & \prob IX $=\ansblank$ \\[12pt]
\prob VI $=\ansblank$ & \prob VIII $=\ansblank$ & \prob X $=\ansblank$ & \prob II $=\ansblank$
\end{tabular}

\bigskip\bigskip
\noindent{\color{sd-blue}\textbf{Part II — Write the Roman numeral.}} \hfill \textit{(Problems 9–16)}

\bigskip
\noindent\begin{tabular}{@{}p{3cm}p{3cm}p{3cm}p{3cm}@{}}
\prob $5 =$ \ansblank & \prob $3 =$ \ansblank & \prob $8 =$ \ansblank & \prob $4 =$ \ansblank \\[12pt]
\prob $9 =$ \ansblank & \prob $6 =$ \ansblank & \prob $7 =$ \ansblank & \prob $10 =$ \ansblank
\end{tabular}

\bigskip
\begin{checkbox}
\textbf{\color{sd-green}Roman Numerals Around Us}\\[3pt]
Look for Roman numerals on clock faces, book chapters, and old buildings.\\
\textit{Can you find a clock with Roman numerals at home or in your neighbourhood?}\\[3pt]
\textbf{Curiosity question:} Why is 4 written as IV (``one before five'') and not IIII?\\
\textit{(Both were used in history — IIII appears on many old clock faces even today!)}
\end{checkbox}

\end{document}
